
% feina a fer per al curs 2010-2011
% 1) l'exercici proposat de maximitzar una funció de dues variables és difícil i cal treballar-lo. Mirar bé el mètode d'steepest descent i el de Newton Raphson (els dos proposats a l'exercici). Mirar també d'entendre correctament el mètode d'optimitzar la funció de Davidson. Mirar una bona explicació a http://linneus20.ethz.ch:8080/1_5_3.html#SECTION00253100000000000000 Fer un dibuis que mostri el concepte de "constant norm", entès com un radi determinat al voltant d'x, un radi donat per l'stepsize. És prou entendor així
\usepackage{framed}
\usepackage{url}
\usepackage{ifthen}
\usepackage{longtable}
\usepackage{fancyvrb}
\usepackage[catalan]{babel}
\usepackage{cancel}
\usepackage{enumitem}
\usepackage[thinc]{esdiff}

%\usepackage[utf8]{inputenc}
\usepackage[catalan]{babel}
\usepackage{lmodern}
\usepackage{diagbox}
\usepackage{amsmath,amsthm,amsfonts,amssymb,amscd}
%\newcommand*\diff{\mathop{}\!\mathrm{d}}
%\newcommand*\Diff[1]{\mathop{}\!\mathrm{d^#1}}
\usepackage[]{systeme}
\usepackage{caption}
\usepackage[arrowdel]{physics}
\usepackage{multirow,booktabs}
\usepackage[dvipsnames,table]{xcolor}
%\usepackage{fullpage}
\usepackage{lastpage}
\usepackage{graphicx}
\graphicspath{{../figures/}}

\usepackage{enumitem}
\usepackage{fancyhdr}
\setlength{\headheight}{14pt}
\usepackage{mathrsfs}
\usepackage{wrapfig}
\usepackage{setspace}
\usepackage{calc}
\usepackage{multicol}
\usepackage{textcomp}
\usepackage{gensymb}
\usepackage{listings}
\usepackage{siunitx}

\usepackage{cancel}
%\usepackage[retainorgcmds]{IEEEtrantools}
\usepackage[margin=3cm]{geometry}
\usepackage{amsmath}
\newlength{\tabcont}
\setlength{\parindent}{0.0in}
\setlength{\parskip}{0.05in}
\usepackage{empheq}
\usepackage{mdframed}

\usepackage[most]{tcolorbox}
% \usepackage{chemfig,chemmacros,chemnum}
% \usepackage{chemformula}

%\usepackage{unicode-math}

\usepackage{tcolorbox}
\usepackage{url}
  \let\oldurl\url
\usepackage{hyperref}
  \let\linkurl\url
  \let\url\oldurl

% \usepackage[
% backend=biber,
% citestyle=numeric-comp
% ]{biblatex}
% \addbibresource{/Users/556079/Documents/teaching.bib}


%\chemsetup[chemformula]{format=\sffamily}
%\renewcommand*\printatom[1]{\ensuremath{\mathsf{#1}}}
%\setatomsep{2em}
%\setdoublesep{.6ex}
%\setbondstyle{semithick}
%\colorlet{shadecolor}{orange!15}
% \parindent 0in
% \parskip 12pt
% \geometry{margin=1in, headsep=0.25in}
% \theoremstyle{definition}
% \newtheorem{defn}{Definition}
% \newtheorem{reg}{Rule}
% \newtheorem{exer}{Exercise}
% \newtheorem{note}{Note}
%\RequirePackage{mathrsfs}
%\RequirePackage[psamsfonts]{amsfonts} %for Y&Y BSR AMS fonts
\RequirePackage{amsmath,amsfonts,amsthm,amssymb}
\RequirePackage{setspace}
\RequirePackage{fancyhdr}
\RequirePackage{lastpage}
\RequirePackage{extramarks}
\RequirePackage{chngpage}
\RequirePackage{soul}


\usepackage{graphicx}
\usepackage{multicol}
\usepackage{hyperref}
\usepackage{makecell}
\usepackage{fancybox}

%\RequirePackage[dvipsnames]{xcolor}
%\RequirePackage{graphicx,float,wrapfig}
\RequirePackage{pgf,tikz}
\usepackage{pgfplots}
%\usetikzlibrary{arrows,automata}
\usepackage{pstricks}
%\RequirePackage{pstricks}
%\RequirePackage[text]{amsthm}
%\RequirePackage{array}
%\RequirePackage{amscd}
%\RequirePackage{array}\RequirePackage{dcolumn}
%\putfig{3.5truein}{PSfig1.3}{Peter's winnings in 40 plays of heads or tails.}{fig 1.3}

% \newcommand*{\DIRFig}{../../../../talks/figures}
% \newcommand{\barefig}[2]{\makebox{\includegraphics*[width=#2] {\DIRFig/#1}}}
% \newcommand{\putfig}[4]
% {\begin{figure}
% \centerline{\barefig{#2}{#1}}
% \caption{#3}
% \label{#4}
% \end{figure}}

\newcommand{\dif}{\, \mathop{}\!\mathrm{d}}
\newcommand{\dx}{\, dx}
\newcommand{\dy}{\, dy}
\newcommand{\dt}{\, dt}
\newcommand{\dth}{\, d\theta}
\newcommand{\dr}{\, dr}
\newcommand{\du}{\, du}
\newcommand\uvec[1]{\textbf{#1}}
\newcommand{\iu}{\hat{\uvec{i}}}
\newcommand{\ju}{\hat{\uvec{j}}}
\newcommand{\ku}{\hat{\uvec{k}}}

% \newcommand{\emx}[1]{{\em{#1}\/}}
% \newcommand{\abin}{{\it ab initio}}
% \newcommand{\bs}{\boldsymbol}
% \newcommand{\citenum}{\cite}
% \newcommand{\dGo}{\ensuremath{\Delta G_0}}
% \newcommand{\dG}[2]{\ensuremath{\Delta G_{\rm #1}^{\rm #2}}}
% \newcommand{\dX}[3]{\ensuremath{\Delta #1_{\rm #2}^{\rm #3}}}
% \newcommand{\ddgo}[1]{\ensuremath{\Delta \Delta G_{\rm solv}^{\rm #1}}}
% \newcommand{\ddgstarcat}{\ensuremath{\Delta \Delta g^{\ddagger}_{\rm cat}}}
% \newcommand{\ddgstar}{\ensuremath{\Delta \dgstar}}
% \newcommand{\ddgt}[2]{\ensuremath{\Delta \Delta G_{\rm solv}^{\rm #1, \rm #2}}}
% \newcommand{\ddsstarprime}{\ensuremath{(\Delta \dsstar)'}}
% \newcommand{\deltaepsel}{\ensuremath{\Delta \varepsilon_{\rm el}}}
% \newcommand{\deltaeps}{\ensuremath{\Delta \varepsilon}}
% \newcommand{\dgab}[2]{\ensuremath{\Delta g_{\rm #1}^{\rm #2}}}
% \newcommand{\dga}[1]{\ensuremath{\Delta g_{\rm #1}}}
% \newcommand{\dgb}[1]{\ensuremath{\Delta g^{\rm #1}}}
% \newcommand{\dgcage}{\ensuremath{\Delta g_{\rm cage}}}
% \newcommand{\dgcat}{\ensuremath{\Delta g_{\rm cat}}}
% \newcommand{\dgsoltsatsa}{\ensuremath{\dgsol (\rm TSA)_{\rm TSA}}}
% \newcommand{\dgsoltstsa}{\ensuremath{\dgsol (\rm TS)_{\rm TSA}}}
% \newcommand{\dgsoltsts}{\ensuremath{\dgsol (\rm TS)_{\rm TS}}}
% \newcommand{\dgsol}{\ensuremath{\Delta G_{\rm sol}}}
% \newcommand{\dgstarcage}{\ensuremath{\dgstar_{\rm cage}}}
% \newcommand{\dgstarcat}{\ensuremath{\dgstar_{\rm cat}}}
% \newcommand{\dgstarw}{\ensuremath{\dgstar_{\rm w}}}
% \newcommand{\dgstar}{\ensuremath{\Delta g^{\ddagger}}}
% \newcommand{\dgw}{\ensuremath{\Delta g_{\rm w}}}
% \newcommand{\dg}[2]{\ensuremath{\Delta g_{\rm #1}^{\rm #2}}}
% \newcommand{\dino}{\texttt{DINO}}
% \newcommand{\dsstarcageprime}{\ensuremath{(\dsstarcage)'}}
% \newcommand{\dsstarcage}{\ensuremath{\dsstar_{\rm cage}}}
% \newcommand{\dsstarcatprime}{\ensuremath{(\dsstarcat)'}}
% \newcommand{\dsstarcat}{\ensuremath{\dsstar_{\rm cat}}}
% \newcommand{\dsstarwprime}{\ensuremath{(\dsstarw)'}}
% \newcommand{\dsstarw}{\ensuremath{\dsstar_{\rm w}}}
% \newcommand{\dsstar}{\ensuremath{\Delta S^{\ddagger}}}
% \newcommand{\eg}{{\it e.g.}}
% \newcommand{\etal}{{\it et al.}}
% \newcommand{\gamess}{\texttt{GAMESS}}
% \newcommand{\gauss}{\texttt{GAUSSIAN} 98}
% \newcommand{\golpe}{\texttt{GOLPE}}
% \newcommand{\grid}{\texttt{GRID}}
% \newcommand{\ie}{{\it i.e.}}
% \newcommand{\ith}{{\it i}$^{\rm th}$\ }
% \newcommand{\kbt}{\ensuremath{k_{\rm B} T}}
% \newcommand{\kb}{\ensuremath{k_{\rm B}}}
% \newcommand{\kcage}{\ensuremath{k_{\rm cage}}}
% \newcommand{\kcatkm}{\ensuremath{k_{\rm cat}/K_{\rm M}}}
% \newcommand{\kcat}{\ensuremath{k_{\rm cat}}}
% \newcommand{\km}{kcal mol$^-1$}
% \newcommand{\knon}{\ensuremath{k_{\rm non}}}
% \newcommand{\kw}{\ensuremath{k_{\rm w}}}
% \newcommand{\mepsim}{\texttt{MEPSIM}}
% \newcommand{\mgp}[1]{\marginpar{\scriptsize{#1}}}
% \newcommand{\mipsim}{\texttt{MIPSIM}}
% \newcommand{\mola}{\texttt{MOLARIS}}
% \newcommand{\msms}{\texttt{MSMS}}
% \newcommand{\pdras}{p21$^{\rm ras}$}
\newcommand{\rgran}{\ensuremath{\mathbb{R}}}
\newcommand{\rx}[2]{\ensuremath{#1_{\rm #2}}}
\newcommand{\vs}{{\it vs.}}
\newcommand{\z}[1]{\ensuremath{\mathbf{#1}}}
\newcommand{\composed}[2]{#1\mathbin\circ #2}
\newcommand{\wrt}[1]{{\mbox{\scriptsize w.r.t. \( #1 \)} }}
\newcommand{\polyspace}{\mathcal{P}}
\newcommand{\matspace}{\mathcal{M}}
%\newcommand{\C}{\mathbb{C}}
\newcommand{\N}{\mathbb{N}}
\newcommand{\Q}{\mathbb{Q}}
\newcommand{\Z}{\mathbb{Z}}
\renewcommand{\Re}{\mathbb{R}}
\newcommand{\rtres}{\ensuremath{\Re^3}}
\newcommand{\union}{\cup}
\newcommand{\dotprod}{\cdot}
\newcommand*\pkg[1]{\textsf{#1}}

%\def\checkmark{\tikz\fill[scale=0.4](0,.35) -- (.25,0) -- (1,.7) -- (.25,.15) -- cycle;}


\newcommand{\trans}[1]{{#1}^{\ensuremath{\mathsf{T}}}} % transpose
\newcommand{\nbyn}[1]{\ensuremath{#1 \! \times \! #1 }}
\newcommand{\nbym}[2]{#1 \! \times \! #2 }       % Use in math mode.
\newcommand{\cat}[2]{#1\!\mathbin{\raise.6ex\hbox{\( {}^\frown \)}}\!#2}
\newcommand{\generalmatrix}[3]{ %arg1: low-case letter, arg2: rows, arg3: cols
               \left(
                  \begin{array}{cccc}
                    #1_{1,1}  &#1_{1,2}  &\ldots  &#1_{1,#2}  \\
                    #1_{2,1}  &#1_{2,2}  &\ldots  &#1_{2,#2}  \\
                              &\vdots                         \\
                    #1_{#3,1} &#1_{#3,2} &\ldots  &#1_{#3,#2}
                  \end{array}
               \right)  }
\newcommand{\colvec}[1]{\begin{pmatrix} #1 \end{pmatrix}}
\newcommand{\pr}[1]{\ensuremath{\mathrm{Pr}(#1)}}
\newcommand{\rep}[2]{ {\rm Rep}_{#2}(#1) }
\newcommand{\mapsunder}[1]{\stackrel{#1}{\longmapsto}}
\newcommand{\map}[3]{\mbox{$#1\colon #2\to #3$}}
\newcommand{\identity}{\mbox{id}}
\newcommand{\stdbasis}{{\cal E}}
\newcommand{\sequence}[1]{ \langle#1\rangle }
\newcommand{\spacer}{\rule[-3mm]{0mm}{8mm}}
\newcommand{\email}[1]{\url{#1}}
\newcommand{\zero}{\vec{0}}
\newcommand{\proj}[2]{\mbox{proj}_{#2}({#1}) }
%\AtBeginDocument{\newlength{\heightofcdot}
%\newlength{\widthofcdot}
%\settoheight{\heightofcdot}{$\cdot$}
%\settowidth{\widthofcdot}{$\cdot$}
%\newsavebox{\dotprodcircle}
%\savebox{\dotprodcircle}{\includegraphics{dotprod.1}}
%\newcommand{\dotprod}{\mathbin{\raisebox{.5\heightofcdot}{%
%          \makebox[\widthofcdot]{$\smash{\usebox{\dotprodcircle}}$}}}}}
\newcommand{\spanof}[1]{\relax [#1\relax ]} % no optional argument!
\newcommand{\set}[1]{\mbox{$\{#1\}$}} \newcommand{\suchthat}{\bigm|}
\newcommand{\deter}[1]{ \mathchoice{\left|#1\right|}{|#1|}{|#1|}{|#1|} }
\newcommand{\secuence}[1]{ \langle#1\rangle }
\newcommand{\basis}[2]{\secuence{\vec{#1}_1,\ldots,\vec{#1}_{#2}}}



%--------linsys
%  Use as \begin{linsys}{3}
%           x &+ &3y &+ &a &= &7 \\
%           x &- &3y &+ &a &= &7
%         \end{linsys}
% Remark: TeXbook pp. 167-170 says to put a medmuskip around a +; and that's
% 4/18-ths of an em.  Why does 2/18-ths of an em work?  I don't know, but
% comparing to a regular displayed equation suggests it is right.
% (darseneau says LaTeX puts in half an \arraycolsep.)
\newenvironment{linsys}[2][m]{%
\setlength{\arraycolsep}{.1111em} % p. 170 TeXbook; a medmuskip
\begin{array}[#1]{@{}*{#2}{rc}r@{}}
}{%
\end{array}}


%\newtheorem{teorema}{Teorema}
%\newtheorem{exercici}{Exercici}
%\newtheorem{definicio}{Definici\'o}
%\newtheorem{theorem}{Theorem}
\newtheorem{exercise}{Exercise}
%\newtheorem{definition}{Definition}

\newcounter{EXMP}
\newenvironment{EXMP}[1][]{\definecolor{shadecolor}{rgb}{0.6,0.6,0.6}
							\begin{shaded}\refstepcounter{EXMP}\par\medskip\noindent%
   							\textbf{EXAMPLE~\theEXMP. #1} \rmfamily}
   							{\end{shaded}\medskip}

\newcounter{BOXT}
\newenvironment{BOXT}[1][]{\definecolor{shadecolor}{rgb}{0.8,0.8,0.8}
							\begin{shaded}\refstepcounter{BOXT}\par\medskip\noindent%
   							\textbf{BOX~\theBOXT. #1} \rmfamily}
   							{\end{shaded}\medskip}

\parskip 4mm


\usepackage{makeidx}
\makeindex


% margins
\topmargin=-0.45in      %
\evensidemargin=0in     %
\oddsidemargin=0in      %
\textwidth=6in        %
\textheight=8.5in       %
\headsep=0.25in         %

% header and footer
\pagestyle{fancy}       %                %
\cfoot{\includegraphics[width=4.8cm]{FCTE}}                %
\rfoot{\thepage}        %
\renewcommand\headrulewidth{0.4pt}   %
\renewcommand\footrulewidth{0.4pt}   %

%\setcounter{section}{-1}

\theoremstyle{definition}
\newtheorem{thm}{Theorem}
\newtheorem{dfn}{Definition}
\newtheorem{lem}{Lemma}
\newtheorem{prp}{Proposition}



%%%%%%%%%%%%%%%%%%%%%%%%%%%%%%%%%%%%%%%%%
%%%%%%%%%%%%%%%%%%%%%%%%%%%%%%%%%%%%%%%%%
% lecturer or student text
% in principle the lecturer text includes some examples to be done in the c lass
\newboolean{LECT}
\setboolean{LECT}{false}
\setboolean{LECT}{true}
%%%%%%%%%%%%%%%%%%%%%%%%%%%%%%%%%%%%%%%%%
%%%%%%%%%%%%%%%%%%%%%%%%%%%%%%%%%%%%%%%%%

\newenvironment{lect}{ %
	\definecolor{shadecolor}{rgb}{1.0,0.8,0.8} %
	\begin{shaded} %
	\textcolor{BrickRed}{\bf Resultat\\}%

} %
{ %
	\end{shaded}
} %

\newcommand{\lct}[1]{\ifthenelse{\boolean{LECT}}{\begin{lect} #1 \end{lect}}{}}



%%%%%%%%%%%%%%%%%%%
% ANGLÈS
%%%%%%%%%%%%%%%%%%%

%\newcommand{\problemName}{}%
%\newcounter{problemCounter}%
%\newenvironment{problem}[1][Problem \arabic{problemCounter}]%
%	{\stepcounter{problemCounter}%
%		\renewcommand{\problemName}{#1}%
%		\section*{\problemName}%
%		\nobreak\extramarks{\problemName}{\problemName continued on next page\ldots}\nobreak%
%		\nobreak\extramarks{\problemName (continued)}{\problemName continued on next page\ldots}\nobreak}%
%	{\nobreak\extramarks{\problemName (continued)}{\problemName continued on next page\ldots}\nobreak%
%		\nobreak\extramarks{\problemName}{}\nobreak}%

\newenvironment{example}{ %merges with the question content
} %


\newenvironment{introduction}{ %
	\definecolor{shadecolor}{rgb}{1.0,1.0,0.8} %
	\begin{shaded} %
	% \textcolor{BrickRed}{\bf Introduction\\}%
} %
{ %
	\end{shaded}
} %


%%%%%%%%%%%%%%%%%%%
% CATALÀ
%%%%%%%%%%%%%%%%%%%
\newtheorem{teorema}{theorem}
\newenvironment{definicio}{ %
	\definecolor{shadecolor}{rgb}{0.9,1.0,0.8} %
	\begin{shaded} %
	\textcolor{OliveGreen}{\bf Definicio\\}%
} %
{ %
	\end{shaded}
} %

%veure http://en.wikibooks.org/wiki/LaTeX/Advanced_Topics
\newcounter{myc}
%environment for exercises in the class notes
\newenvironment{exr}{ %
    \addtocounter{myc}{1}
	\definecolor{shadecolor}{rgb}{0.9,1.0,0.8} %
	\begin{shaded} %
	\textcolor{OliveGreen}{\bf Exercici \arabic{myc}\\}%
} %
{ %
	\end{shaded}
} %
%environment for questions in exams
\newenvironment{qst}{ %
    \addtocounter{myc}{1}
	\definecolor{shadecolor}{rgb}{0.9,1.0,0.8} %
	\begin{shaded} %
	\textcolor{OliveGreen}{\bf Qüestió \arabic{myc}\\}%
} %
{ %
	\end{shaded}
} %


\usepackage[lastexercise]{exercise}
% \renewcommand{\listexercisename}{\'{I}ndex d'exercicis}
% \renewcommand{\ExerciseListName}{Ex.}  % titol a cada exercici
% \renewcommand{\AnswerListName}{Soluci\'{o}}
% \renewcommand{\AnswerName}{Soluci\'{o} de l'exercici}
% \renewcommand{\ExerciseName}{Exercici}  % títol que apareix al llistat d'exercicis
% \renewcommand{\ExerciseHeaderTitle}{\ExerciseTitle \quad---\quad \medskip}
% \renewcommand{\ExerciseListHeader}{
%         \shadowbox*{
%                 \ExerciseHeaderDifficulty%
%                 \textbf{\ExerciseListName\ \ExerciseHeaderNB%
%                 \ --- \ \ExerciseHeaderTitle}%
%                 \medskip\ExerciseHeaderOrigin%
%                 \ignorespaces
%         }
% }

\makeatletter
% remove the faulty \expandafter
\patchcmd{\@@@ExeEnv}
  {\theExercise\ \expandafter}
  {\theExercise\ }
  {}{}
\patchcmd{\@@@ExeCmd}
  {\theExercise\ \expandafter}
  {\theExercise\ }
  {}{}
% remove \itshape
\patchcmd{\@@@ExeEnv}
  {\itshape}
  {}
  {}{}
\patchcmd{\@@@ExeCmd}
  {\itshape}
  {}
  {}{}
\makeatother

% comanda per definir Gauss-Jordan Elimination
\allowdisplaybreaks
\makeatletter
\newcounter{elimination@steps}
\newcolumntype{R}[1]{>{\raggedleft\arraybackslash$}p{#1}<{$}}
\def\elimination@num@rights{}
\def\elimination@num@variables{}
\def\elimination@col@width{}
\newenvironment{elimination}[4][0]
{
    \setcounter{elimination@steps}{0}
    \def\elimination@num@rights{#1}
    \def\elimination@num@variables{#2}
    \def\elimination@col@width{#3}
    \renewcommand{\arraystretch}{#4}
    \start@align\@ne\st@rredtrue\m@ne
}
{
    \endalign
    \ignorespacesafterend
}
\newcommand{\eliminationstep}[2]
{
    \ifnum\value{elimination@steps}>0\sim\quad\fi
    \left[
        \ifnum\elimination@num@rights>0
            \begin{array}
            {@{}*{\elimination@num@variables}{R{\elimination@col@width}}
            |@{}*{\elimination@num@rights}{R{\elimination@col@width}}}
        \else
            \begin{array}
            {@{}*{\elimination@num@variables}{R{\elimination@col@width}}}
        \fi
            #1
        \end{array}
    \right]
    &
    \begin{array}{l}
        #2
    \end{array}
    &%                                    moved second & here
    \addtocounter{elimination@steps}{1}
}
\makeatother

\usepackage{tkz-euclide}
